% =============================================================
% บทที่ 3 - วิธีการดำเนินงาน
% โครงงาน: การออกแบบและพัฒนาเกมเดินเรือด้วยระบบสั่งการด้วยเสียงเพื่อส่งเสริมทักษะการสื่อสาร
%          (Design and Development of a Voice-Controlled Ship Navigation Game to Enhance Communication Skills)
% =============================================================

% -----------------------------------------------------------
\section{ขั้นตอนการดำเนินงาน}
% -----------------------------------------------------------

การดำเนินงานวิจัยและการพัฒนาเกมเดินเรือด้วยระบบสั่งการด้วยเสียง แบ่งออกเป็น 7 ขั้นตอนหลัก ดังนี้

\subsection{การวิเคราะห์ปัญหาและศึกษาทฤษฎีที่เกี่ยวข้อง}

\begin{enumerate}[label=(\arabic*)]
    \item ศึกษาค้นคว้าทฤษฎีที่เกี่ยวข้องกับการเรียนรู้โดยใช้เกมเป็นฐาน (DGBL) สมมติฐานตัวกรองทางอารมณ์ (Affective Filter Hypothesis) ของ Krashen และทฤษฎีการเรียนรู้ตามสถานการณ์ (Situated Learning) เพื่อนำมาเป็นกรอบแนวคิดในการออกแบบสภาพแวดล้อมที่ช่วยลดความวิตกกังวลในการพูดภาษาอังกฤษ
    \item ศึกษาคำศัพท์และรูปแบบประโยคมาตรฐานที่ใช้ในการสื่อสารเพื่อการเดินเรือ (Standard Marine Communication Phrases) เช่น การสั่งหันหางเสือ (Port/Starboard) และการสั่งระดับความเร็ว เพื่อนำมาสร้างเป็นคลังคำสั่ง (Command Vocabulary) ภายในเกม
    \item ศึกษาเทคโนโลยีและเครื่องมือสำหรับการพัฒนา ได้แก่ Unity ภาษา C\# รวมถึงโมเดลการรู้จำเสียง ggml-base.en (ผ่านไลบรารี whisper.cpp) และแบบจำลองภาษา GPT-4o mini ผ่าน OpenAI API
\end{enumerate}

\subsection{การกำหนดขอบเขตและข้อกำหนดของระบบ}

กำหนดขอบเขตของเกมให้เป็นการจำลองสถานการณ์บังคับเรือหลบหลีกสิ่งกีดขวาง โดยใช้เสียงพูดภาษาอังกฤษผ่านไมโครโฟนเป็นวิธีการหลักในการสั่งการทิศทางและความเร็วของเรือ ผู้เล่นพิมพ์คำสั่งเป็นข้อความได้เป็นช่องทางสำรอง และไม่มีปุ่มใดที่ควบคุมหางเสือหรือเครื่องยนต์โดยตรง ทุกคำสั่งต้องสื่อสารผ่านลูกเรือ และกำหนดกลุ่มตัวอย่างสำหรับการประเมินผลจำนวน 10--20 คน

\subsection{การออกแบบระบบและเกม (Game Design)}

\begin{enumerate}[label=(\arabic*)]
    \item \textbf{การออกแบบเนื้อเรื่องและบทบาท:} ผู้เล่นสวมบทบาทเป็นกัปตันเรือที่ต้องสั่งการลูกเรือปัญญาประดิษฐ์ (AI Crew) ได้แก่ ต้นกลเรือ และพลถือท้ายและยามระวังหน้า เพื่อนำเรือไปถึงจุดหมายอย่างปลอดภัย
    \item \textbf{การออกแบบกลไกหลัก (Core Mechanics):} ผู้เล่นกดปุ่มวิทยุค้างไว้แล้วพูดคำสั่ง เมื่อปล่อยปุ่ม ระบบแปลงเสียงเป็นข้อความและส่งให้ตัวแยกคำสั่ง (Command Parser) ตีความ หากพบคำสั่งที่ชัดเจน ลูกเรือจะปฏิบัติตามและเรือเคลื่อนที่ตามหลักฟิสิกส์
    \item \textbf{การออกแบบระบบตอบกลับ (Feedback System):} ออกแบบให้ลูกเรือปัญญาประดิษฐ์ทวนคำสั่ง หรือร้องขอให้ผู้เล่นย้ำคำสั่งใหม่เมื่อการสื่อสารติดขัด
    \item \textbf{การออกแบบ UI/UX:} ออกแบบหน้าจอแดชบอร์ด เข็มทิศ เรดาร์ตรวจจับสิ่งกีดขวาง และแถบแสดงสถานะการรับเสียงของไมโครโฟน
\end{enumerate}

\subsection{การพัฒนาเกม (Game Development)}

ดำเนินการพัฒนาเกมโดยแบ่งออกเป็น 4 ระยะ ได้แก่

\begin{enumerate}[label=(\arabic*)]
    \item \textbf{ระยะที่ 1 (Prototype):} พัฒนาระบบพื้นฐาน เชื่อมต่อไมโครโฟนเข้ากับระบบ Speech-to-Text และทดสอบการแปลงเสียงเป็นคำสั่งเพื่อบังคับการเคลื่อนที่ของเรือเบื้องต้น
    \item \textbf{ระยะที่ 2 (Alpha):} พัฒนาสภาพแวดล้อมทางทะเล ระบบสุ่มสร้างสิ่งกีดขวาง และระบบเสียงโต้ตอบของลูกเรือปัญญาประดิษฐ์
    \item \textbf{ระยะที่ 3 (Beta):} พัฒนาระบบ UI/UX ให้สมบูรณ์ ปรับจูนระดับความยาก และปรับแต่งความแม่นยำในการรับคำสั่งเสียง
    \item \textbf{ระยะที่ 4 (Release Candidate):} ตรวจสอบและแก้ไขข้อผิดพลาด (Bug Fixing) และจัดเตรียมแพ็กเกจตัวเกมสำหรับนำไปทดสอบกับกลุ่มตัวอย่าง
\end{enumerate}

\subsection{การทดสอบระบบ (Testing)}

ดำเนินการทดสอบระบบเป็น 4 ช่วง ดังนี้

\begin{enumerate}[label=(\arabic*)]
    \item \textbf{ช่วงที่ 1 (Unit Testing):} ทดสอบตรรกะของตัวแยกคำสั่ง ระบบช่วยสอนภาษา และระบบความก้าวหน้าของผู้เล่น (อู่ต่อเรือ อัปเกรด และโหมดเนื้อเรื่อง) ด้วย Unity Test Framework โดยไม่ต้องใช้ไมโครโฟนหรือเปิดฉากเกม
    \item \textbf{ช่วงที่ 2 (Integration and System Testing):} ทดสอบการทำงานร่วมกันตั้งแต่การรับเสียง การแปลงเสียงเป็นข้อความ ตัวแยกคำสั่ง การตอบกลับของลูกเรือ ระบบฟิสิกส์เรือ และระบบสิ่งกีดขวาง ตามกรณีทดสอบที่กำหนด รวมถึงกรณีที่ผู้เล่นออกคำสั่งที่ไม่คาดคิด
    \item \textbf{ช่วงที่ 3 (Non-functional Testing):} วัดเวลาตอบสนอง (Latency) ของการแปลงเสียงเป็นข้อความและการตอบสนองของเรือจากแฟ้มบันทึกผล และวัดความแม่นยำของระบบจากคำสั่งเป้าหมายที่ผู้เข้าร่วมอ่าน
    \item \textbf{ช่วงที่ 4 (User Acceptance Testing: UAT):} ให้กลุ่มตัวอย่างทดลองเล่นตามขั้นตอนการประเมินผล เพื่อตรวจสอบความเข้าใจในระบบและค้นหาจุดบกพร่อง
\end{enumerate}

\subsection{การประเมินผล (Evaluation)}

การประเมินผลใช้แบบแผนการทดลองกลุ่มเดียววัดก่อนและหลัง (One-Group Pretest--Posttest Design) กับกลุ่มตัวอย่าง 10--20 คน โดยมีรายละเอียดดังนี้

\begin{enumerate}[label=\arabic*)]
    \item \textbf{ประชากรและกลุ่มตัวอย่าง:} ประชากร ได้แก่ นักศึกษาระดับปริญญาตรี กลุ่มตัวอย่างเลือกแบบอาสาสมัคร (Convenience Sampling) โดยมีเกณฑ์คัดเข้า ได้แก่ เป็นนักศึกษาระดับปริญญาตรี ใช้ภาษาไทยเป็นภาษาแม่ และสามารถใช้ชุดหูฟังที่มีไมโครโฟนได้ จำนวน 10--20 คน ซึ่งเป็นขนาดของการศึกษานำร่อง (Pilot Study) ตามเวลาและทรัพยากรของโครงงาน เมื่อคำนวณด้วย G*Power สำหรับการทดสอบค่าทีแบบกลุ่มตัวอย่างไม่อิสระ ทางเดียว ($\alpha$ = 0.05, power = 0.80) กลุ่มตัวอย่างขนาดนี้ตรวจพบได้เฉพาะผลขนาดใหญ่ (d ประมาณ 0.8 ซึ่งต้องการอย่างน้อย 12 คน) จึงรายงานขนาดอิทธิพลควบคู่กับค่านัยสำคัญเสมอ และหากผู้เข้าร่วมที่ตอบครบทั้งก่อนและหลังมีน้อยกว่า 12 คน จะรายงานผลเชิงพรรณนาเป็นหลัก
    \item \textbf{สถานที่และอุปกรณ์:} เก็บข้อมูลรายบุคคลในห้องที่เงียบ ผู้เข้าร่วมทุกคนใช้เครื่องคอมพิวเตอร์ทดสอบเครื่องเดียวกันและชุดหูฟังที่มีไมโครโฟนตามตารางที่ \ref{tab:hardware}
    \item \textbf{การขอความยินยอม:} ผู้เข้าร่วมทุกคนให้ความยินยอมก่อนเริ่มการทดลอง เอกสารขอความยินยอมระบุว่าเกมไม่บันทึกเสียงของผู้เข้าร่วม แต่เกมอาจส่งข้อความที่ถอดจากเสียงไปประมวลผลผ่าน OpenAI API และแบบสอบถามไม่เก็บชื่อหรือข้อมูลที่ระบุตัวตนได้
    \item \textbf{เครื่องมือ:} (1) ตอนที่ 1 ข้อมูลทั่วไป 6 ข้อ (2) ตอนที่ 2 แบบวัดความวิตกกังวลในการพูดภาษาอังกฤษ 12 ข้อ แบ่งเป็นด้านความประหม่าในการสื่อสาร (Communication Apprehension: CA) 6 ข้อ และด้านความกลัวการถูกประเมินในทางลบ (Fear of Negative Evaluation: FNE) 6 ข้อ มีข้อความเชิงบวก 4 ข้อ ข้อความทุกข้อเขียนขึ้นใหม่เป็นภาษาไทยตามองค์ประกอบของ FLCAS~\cite{horwitz1986} ไม่ใช่การแปลข้อความต้นฉบับ (3) ตอนที่ 3 แบบสอบถามประสบการณ์การเล่นเกม 10 ข้อ (4) ตอนที่ 4 คำถามปลายเปิด 3 ข้อ (ภาคผนวก ก) (5) รายการคำสั่งเป้าหมาย 10 คำสั่ง และ (6) แฟ้มบันทึกผลการเล่นจากเกม (turns.csv, voyages.csv และ emergencies.csv) ซึ่งแยกผู้เข้าร่วมแต่ละคนด้วยรอบการเปิดเกม (คอลัมน์ \texttt{session}) โดยไม่เชื่อมกับคำตอบในแบบสอบถาม
    \item \textbf{การตรวจสอบคุณภาพเครื่องมือ:} ตรวจความตรงเชิงเนื้อหาด้วยดัชนีความสอดคล้อง (IOC) จากผู้เชี่ยวชาญ 3 ท่าน โดยคงไว้เฉพาะข้อที่มีค่า IOC ตั้งแต่ 0.50 ขึ้นไป จากนั้นนำไปทดลองใช้ (Try-out) กับนักศึกษาที่ไม่ใช่กลุ่มตัวอย่าง 30 คน และหาความเชื่อมั่นด้วยสัมประสิทธิ์แอลฟาของครอนบาค (Cronbach's Alpha)~\cite{cronbach1951} แยกเป็นด้าน CA ด้าน FNE และตอนที่ 3 โดยเกณฑ์ที่ยอมรับได้คือ 0.70 ขึ้นไป
    \item \textbf{ขั้นตอนการเก็บข้อมูล:} (1) ผู้ดำเนินการทดลองชี้แจงวัตถุประสงค์และขอความยินยอม (2) ผู้เข้าร่วมเปิดแบบสอบถามออนไลน์ซึ่งเป็นฟอร์มเดียวตลอดการทดลอง และตอบตอนที่ 1 และ 2 ก่อนเล่น (3) ผู้ดำเนินการทดลองเปิดเกมใหม่สำหรับผู้เข้าร่วมแต่ละคน เพื่อให้แฟ้มบันทึกผลของแต่ละคนแยกกันด้วยรอบการเปิดเกม (4) ผู้เข้าร่วมเล่นบทฝึกสอนและเดินเรือต่อเนื่องประมาณ 30 นาที โดยเปิดฟอร์มค้างไว้ ระหว่างเล่นผู้เข้าร่วมอ่านรายการคำสั่งเป้าหมาย 10 คำสั่งผ่านวิทยุ (5) ผู้เข้าร่วมกลับมาที่ฟอร์มเดิม ตอบตอนที่ 2 ซ้ำ และตอบตอนที่ 3 และ 4 แล้วจึงส่งฟอร์ม คำตอบก่อนและหลังเล่นของผู้เข้าร่วมแต่ละคนจึงอยู่ในการส่งครั้งเดียวกัน
    \item \textbf{การวิเคราะห์ข้อมูล:} กลับคะแนนข้อที่มีเครื่องหมาย (R) ในตอนที่ 2 และข้อ 10 ในตอนที่ 3 (6 $-$ คะแนนที่ตอบ) แล้วใช้ค่าเฉลี่ยรายข้อของผู้เข้าร่วมแต่ละคน (ช่วง 1--5) เป็นหน่วยวิเคราะห์ หาค่าเฉลี่ยและส่วนเบี่ยงเบนมาตรฐาน และแปลความหมายตามเกณฑ์ในตารางที่ \ref{tab:likert_criteria} สำหรับการเปรียบเทียบความวิตกกังวลก่อนและหลังเล่น ตรวจการแจกแจงของผลต่างคะแนนด้วยการทดสอบชาปิโร--วิลค์ (Shapiro--Wilk Test) หากแจกแจงแบบปกติใช้การทดสอบค่าที (t-test) แบบกลุ่มตัวอย่างไม่อิสระ (Paired-Samples t-test) และรายงานขนาดอิทธิพล (Effect Size) Cohen's $d_z$ หากไม่แจกแจงแบบปกติใช้การทดสอบวิลคอกซัน (Wilcoxon Signed-Rank Test) และรายงานขนาดอิทธิพล $r = Z/\sqrt{N}$ โดยทดสอบแบบทางเดียวเนื่องจากสมมติฐานระบุทิศทาง ที่ระดับนัยสำคัญ 0.05 สำหรับวัตถุประสงค์ด้านประสบการณ์การเล่นเกม กำหนดเกณฑ์ผ่านที่ค่าเฉลี่ยประสบการณ์การเล่นเกมตั้งแต่ 3.41 ขึ้นไป (ระดับมาก) ส่วนประสิทธิภาพทางเทคนิควัดจากแฟ้มบันทึกผล ได้แก่ ความแม่นยำซึ่งคิดเป็นร้อยละของคำสั่งเป้าหมายที่ระบบถอดความและแยกคำสั่งได้ถูกต้อง โดยเทียบรายการคำสั่งเป้าหมายกับคอลัมน์ข้อความที่ได้ยินและคำสั่งที่ปฏิบัติใน turns.csv ร้อยละของการสื่อสารที่ลูกเรือเข้าใจระหว่างการเล่นอิสระ จำนวนครั้งที่ลูกเรือขอให้พูดซ้ำ และเวลาตอบสนองของระบบ
    \item \textbf{สมมติฐาน:} คะแนนความวิตกกังวลในการพูดภาษาอังกฤษของผู้เล่นหลังเล่นเกมต่ำกว่าก่อนเล่นเกมอย่างมีนัยสำคัญทางสถิติที่ระดับ 0.05
\end{enumerate}

แบบแผนกลุ่มเดียววัดก่อนและหลังมีภัยคุกคามต่อความตรงภายในที่ต้องคำนึงถึง ได้แก่ (1) ผลจากการทดสอบ (Testing Effect) เนื่องจากผู้เข้าร่วมตอบข้อความ 12 ข้อชุดเดิมซ้ำหลังผ่านไปประมาณ 30 นาที (2) การถดถอยเข้าหาค่าเฉลี่ย (Regression to the Mean) โดยเฉพาะผู้ที่มีคะแนนก่อนเล่นสูงหรือต่ำมาก (3) ผลจากความแปลกใหม่และการรู้ตัวว่าถูกศึกษา (Novelty and Hawthorne Effect) (4) การตอบตามความคาดหวังของสังคม (Social Desirability) เนื่องจากผู้ดำเนินการทดลองอยู่ในห้องขณะตอบแบบสอบถาม และ (5) การไม่มีกลุ่มควบคุมและไม่มีการวัดซ้ำหลังผ่านไประยะหนึ่ง (Delayed Post-test) ดังนั้นผลการประเมินจะแปลความเป็นความวิตกกังวลในการพูดภาษาอังกฤษระยะสั้นขณะทดลอง (State Anxiety) ภายในการทดลองครั้งนั้นเท่านั้น ไม่ใช่การเปลี่ยนแปลงในระยะยาว และไม่สรุปความเป็นเหตุเป็นผลอย่างเด็ดขาด

เกณฑ์การแปลความหมายค่าเฉลี่ยแบ่งเป็น 5 ระดับ โดยมีความกว้างของอันตรภาคชั้นเท่ากับ (5 $-$ 1)/5 = 0.80 ดังตารางที่ \ref{tab:likert_criteria}

\begin{table}[!htbp]
\centering
\caption{เกณฑ์การแปลความหมายค่าเฉลี่ย}
\label{tab:likert_criteria}
\begin{tabular}{|L{2.91cm}|L{4.85cm}|L{5.52cm}|}
\hline
\textbf{ค่าเฉลี่ย} & \textbf{ระดับความวิตกกังวล (ตอนที่ 2)} & \textbf{ระดับประสบการณ์การเล่นเกม (ตอนที่ 3)} \\
\hline
4.21--5.00 & สูงมาก & มากที่สุด \\ \hline
3.41--4.20 & สูง & มาก \\ \hline
2.61--3.40 & ปานกลาง & ปานกลาง \\ \hline
1.81--2.60 & ต่ำ & น้อย \\ \hline
1.00--1.80 & ต่ำมาก & น้อยที่สุด \\
\hline
\end{tabular}
\end{table}

\subsection{การจัดทำรายงานและสรุปผล}

รวบรวมข้อมูลจากการออกแบบ การพัฒนา ผลการทดสอบ และการประเมินผล มาวิเคราะห์เพื่อจัดทำรูปเล่มรายงานโครงงานฉบับสมบูรณ์

% -----------------------------------------------------------
\section{เครื่องมือที่ใช้ในการดำเนินงาน}
% -----------------------------------------------------------

\subsection{ฮาร์ดแวร์}

รายการฮาร์ดแวร์ที่ใช้ในการพัฒนาและทดสอบโครงงานแสดงดังตารางที่ \ref{tab:hardware}

\begin{table}[!htbp]
\centering
\caption{รายการฮาร์ดแวร์ที่ใช้ในการพัฒนาและทดสอบ}
\label{tab:hardware}
\begin{tabular}{|c|L{3.60cm}|L{8.30cm}|}
\hline
\textbf{ลำดับ} & \textbf{อุปกรณ์} & \textbf{รายละเอียด/คุณสมบัติ} \\
\hline
1 & เครื่องคอมพิวเตอร์ทดสอบ & คอมพิวเตอร์โน้ตบุ๊ก หน่วยประมวลผลกลาง (CPU) Intel Core i5-13420H (รุ่นที่ 13) การ์ดจอ (GPU) NVIDIA GeForce RTX 4050 Laptop GPU และหน่วยความจำ (RAM) 16 GB ใช้รัน Unity และ whisper.cpp \\ \hline
2 & ชุดหูฟังที่มีไมโครโฟน & ชุดหูฟังทั่วไปที่มีไมโครโฟนในตัว ใช้รับเสียงพูดภาษาอังกฤษของผู้เล่น และฟังเสียงสภาพแวดล้อมและเสียงโต้ตอบจากลูกเรือปัญญาประดิษฐ์ \\
\hline
\end{tabular}
\end{table}

\subsection{ซอฟต์แวร์และบริการ}

รายการซอฟต์แวร์และส่วนเชื่อมต่อโปรแกรมประยุกต์ที่ใช้แสดงดังตารางที่ \ref{tab:software}

\begin{table}[!htbp]
\centering
\caption{รายการซอฟต์แวร์และบริการที่ใช้ในการพัฒนา}
\label{tab:software}
\begin{tabular}{|c|L{3.20cm}|L{3.30cm}|L{5.10cm}|}
\hline
\textbf{ลำดับ} & \textbf{ซอฟต์แวร์ / บริการ} & \textbf{เวอร์ชัน/ประเภท} & \textbf{วัตถุประสงค์} \\
\hline
1 & Windows 10/11 & ระบบปฏิบัติการ & ระบบปฏิบัติการหลักสำหรับพัฒนาและรันเกม (เสียงพูดของลูกเรือใช้ระบบสังเคราะห์เสียงของ Windows) \\ \hline
2 & Unity & Unity 6 (6000.5.8f1) & สร้างสภาพแวดล้อม 3 มิติและระบบฟิสิกส์การเดินเรือ \\ \hline
3 & C\# & ภาษาโปรแกรม & ภาษาหลักในการควบคุมกลไกภายในเกม \\ \hline
4 & whisper.cpp & โมเดล ggml-base.en (ทำงานภายในเครื่อง) & แปลงเสียงพูดภาษาอังกฤษเป็นข้อความหลังผู้เล่นปล่อยปุ่มวิทยุ \\ \hline
5 & OpenAI API & GPT-4o mini (บริการคลาวด์) & ตีความข้อความที่ตัวแยกคำสั่งไม่พบคำสั่งที่ชัดเจน และสร้างบทสนทนาของลูกเรือปัญญาประดิษฐ์ \\ \hline
6 & Visual Studio / VS Code & IDE & สภาพแวดล้อมสำหรับการเขียนชุดคำสั่ง \\ \hline
7 & Git / GitHub & Version Control & ระบบควบคุมและจัดการรุ่นของซอร์สโค้ด \\ \hline
8 & Google Forms & บริการออนไลน์ & เก็บแบบสอบถาม \\ \hline
9 & JASP หรือ SPSS & โปรแกรมสถิติ & สำหรับวิเคราะห์ข้อมูล \\
\hline
\end{tabular}
\end{table}

% -----------------------------------------------------------
\section{แผนการดำเนินงาน}
% -----------------------------------------------------------

แผนการดำเนินงานตลอดระยะเวลาการทำโครงงาน ปีการศึกษา 2569 แสดงดังตารางที่ \ref{tab:gantt}

\begin{table}[!htbp]
\centering
\caption{แผนและระยะเวลาดำเนินการโครงงาน}
\label{tab:gantt}
\small
\begin{tabular}{|c|L{4.60cm}|c|c|c|c|c|c|c|c|c|c|}
\hline
\textbf{ลำดับ} & \textbf{กิจกรรม} & \multicolumn{2}{c|}{\textbf{ก.ค. 69}} & \multicolumn{2}{c|}{\textbf{ส.ค. 69}} & \multicolumn{2}{c|}{\textbf{ก.ย. 69}} & \multicolumn{2}{c|}{\textbf{ต.ค. 69}} & \multicolumn{2}{c|}{\textbf{พ.ย. 69}} \\ \cline{3-12}
 & & \textbf{1} & \textbf{2} & \textbf{1} & \textbf{2} & \textbf{1} & \textbf{2} & \textbf{1} & \textbf{2} & \textbf{1} & \textbf{2} \\ \hline
1 & กำหนดหัวข้อโครงงาน & $\bullet$ & & & & & & & & & \\ \hline
2 & ศึกษากลุ่มผู้ใช้และงานวิจัย & $\bullet$ & $\bullet$ & & & & & & & & \\ \hline
3 & กำหนดขอบเขต ข้อจำกัด และเป้าหมาย & & $\bullet$ & & & & & & & & \\ \hline
4 & การเขียนข้อเสนอโครงงาน & & $\bullet$ & $\bullet$ & & & & & & & \\ \hline
5 & นำเสนอข้อเสนอโครงงาน & & & $\bullet$ & & & & & & & \\ \hline
6 & การสร้างและพัฒนาต้นแบบ (Development) & & & $\bullet$ & $\bullet$ & $\bullet$ & $\bullet$ & $\bullet$ & & & \\ \hline
7 & การทดสอบระบบและปรับปรุง & & & & & $\bullet$ & $\bullet$ & $\bullet$ & $\bullet$ & & \\ \hline
8 & สร้างและตรวจคุณภาพเครื่องมือ (IOC, Try-out) & & & & & & $\bullet$ & $\bullet$ & $\bullet$ & & \\ \hline
9 & เก็บข้อมูลและประเมินผลกับกลุ่มตัวอย่าง & & & & & & & & & $\bullet$ & \\ \hline
10 & วิเคราะห์และสรุปผลโครงงาน & & & & & & & & & $\bullet$ & $\bullet$ \\ \hline
11 & จัดทำรูปเล่มโครงงาน & & & & & & $\bullet$ & $\bullet$ & $\bullet$ & $\bullet$ & $\bullet$ \\ \hline
12 & ส่งมอบและนำเสนอโครงงาน & & & & & & & & & & $\bullet$ \\ \hline
\end{tabular}

\vspace{0.2cm}
\raggedright หมายเหตุ: ช่อง 1 หมายถึงครึ่งเดือนแรก และช่อง 2 หมายถึงครึ่งเดือนหลัง การเก็บข้อมูลกับกลุ่มตัวอย่างเริ่มหลังเครื่องมือผ่านการตรวจคุณภาพแล้ว
\end{table}